\documentclass[11pt,a4paper]{article}
\usepackage{fontspec,amsmath,amssymb,amsthm,microtype}
\setmainfont{lmroman10-regular.otf}[BoldFont=lmroman10-bold.otf,ItalicFont=lmroman10-italic.otf,BoldItalicFont=lmroman10-bolditalic.otf]
\setmonofont{lmmono10-regular.otf}[Scale=MatchLowercase]
\usepackage[margin=25mm]{geometry}
\usepackage[unicode,hidelinks]{hyperref}
\setlength{\emergencystretch}{3em}
\setlength{\parskip}{0.4em}
\setlength{\parindent}{0pt}
\allowdisplaybreaks[2]
\raggedbottom
\title{\texorpdfstring{A first Euler-\allowbreak{}step estimate}{A first Euler-step estimate}}
\author{LangevinLean mathematical exposition}
\date{}
\begin{document}\maketitle
\noindent Source module: \texttt{Unadjusted\allowbreak{}Langevin/\allowbreak{}Gradient\allowbreak{}Step.\allowbreak{}lean}. The accompanying HTML pairs this mathematical exposition with the exact Lean statements and records the scope of the formalization.


Before using the full lower and upper models together, it is useful to record the estimate obtained directly from strong monotonicity and a Lipschitz bound. This estimate applies to arbitrary vectors \(g_x,g_y\); they need not arise as gradients.



The formal inputs \texttt{gx} and \texttt{gy} are arbitrary vectors satisfying the displayed monotonicity and Lipschitz inequalities. Thus \texttt{gradient\_\allowbreak{}step\_\allowbreak{}sq\_\allowbreak{}bound} is an algebraic Hilbert-space result; the potential's actual gradients are substituted later. \texttt{gradient\_\allowbreak{}step\_\allowbreak{}contract} takes a nonnegative square root and does not require a separate sign assumption on the scalar coefficient.


\section{\texorpdfstring{The squared Euler-\allowbreak{}step bound}{The squared Euler-step bound}}


Let \(x,y,g_x,g_y\) lie in a real inner-product space, let \(m\in\mathbb R\), \(L\ge0\), and \(h\ge0\). Assume

\[m\|x-y\|^2\le\langle g_x-g_y,x-y\rangle,\qquad\|g_x-g_y\|\le L\|x-y\|.\]

Then

\[\|(x-hg_x)-(y-hg_y)\|^2\le(1-2hm+h^2L^2)\|x-y\|^2.\]

Both sides of the Lipschitz estimate are nonnegative, so squaring preserves the inequality:

\[\|g_x-g_y\|^2\le L^2\|x-y\|^2.\]

Multiply strong monotonicity by \(2h\ge0\), and multiply the squared Lipschitz estimate by \(h^2\ge0\):

\[2hm\|x-y\|^2\le2h\langle g_x-g_y,x-y\rangle,\qquad h^2\|g_x-g_y\|^2\le h^2L^2\|x-y\|^2.\]

Write \(d=x-y\), \(v=g_x-g_y\). Expanding the squared norm and applying the preceding estimates gives

\[\begin{aligned}\|(x-hg_x)-(y-hg_y)\|^2&=\|d-hv\|^2\\&=\|d\|^2-2h\langle v,d\rangle+h^2\|v\|^2\\&\le(1-2hm+h^2L^2)\|d\|^2.\end{aligned}\]

{\raggedright\noindent Lean declarations: \texttt{gradient\_\allowbreak{}step\_\allowbreak{}sq\_\allowbreak{}bound}.\par}

\section{\texorpdfstring{Taking the square root}{Taking the square root}}


Let \(H\) be a real inner-product space, \(x,y,g_x,g_y\in H\), \(m\in\mathbb R\), \(L\ge0\), and \(h\ge0\). Assume

\[m\|x-y\|^2\le\langle g_x-g_y,x-y\rangle,\qquad\|g_x-g_y\|\le L\|x-y\|.\]

Then

\[\|(x-hg_x)-(y-hg_y)\|\le\sqrt{1-2hm+h^2L^2}\,\|x-y\|.\]

No separate sign assumption on \(1-2hm+h^2L^2\) is required; Lean uses the real square root, which is zero on nonpositive inputs.



Apply the preceding squared estimate. With \(A=\|(x-hg_x)-(y-hg_y)\|\), \(r=\|x-y\|\), and \(c=1-2hm+h^2L^2\), it says \(A^2\le cr^2\).



Since \(A\ge0\), take square roots and use \(r^2\ge0\):

\[A\le\sqrt{cr^2}=\sqrt c\,\sqrt{r^2}=\sqrt c\,r.\]

The square-root product identity is valid here even without a sign hypothesis on \(c\), because the second factor \(r^2\) is nonnegative.



{\raggedright\noindent Lean declarations: \texttt{gradient\_\allowbreak{}step\_\allowbreak{}contract}.\par}

\end{document}
